\section{Two prescribed edges and the polygon second moment}\label{c4:metric-consequences}

The analytic estimate is now complete. We identify its normalization
with a positive sum of physical polygons, then sum over displacements
between the two marked edges. Tilted cuts let us treat every displacement
except three lattice directions; those exceptional pairs are controlled
by the first length moment.

\subsection{A tilted cut and its positive two-edge mass}

Use the site order $(P,L,T,R)$ of \eqref{c4:Pair-sum}. Choose a pattern
of $k=a+b$ sites, with $a$ unshifted sites $(z=Q,x=q)$ and $b$ shifted
sites $(z=1,x=1)$, where $1\le b\le a$. Repeat this pattern $m\ge3$
times, with $N=mk$ even, and mark its first and $(k+1)$-st sites.
The marks have the same type, and $L$ is the shorter block between them.

We realize this array at auxiliary parameter $Q^{-1}$ in the ordinary
honeycomb lattice. Tile the plane by rhombi with side vectors $d_j,v$,
where
\[
 h=(\sqrt3,0),\qquad v=(\sqrt3/2,3/2),\qquad
 d_j=\begin{cases}h-v,&z_j=Q,\\h,&z_j=1.\end{cases}
\]
Their lower corners are $\sum_{i<j}d_i+nv$, with $n\in\mathbb Z$ and
the pattern extended in both column directions. Split each rhombus
along its shorter diagonal into two elementary triangles, put a
trivalent vertex at each centroid, and put ports at the outer side
midpoints. At $z_j=1$ the two vertices group bottom-left and right-top
ports; at $z_j=Q$ they group bottom-right and left-top ports. These are
exactly the physical local weights: the resulting dual graph is the
ordinary honeycomb, with weight $x_*$ per visited vertex
\cite{GlazmanManolescu2017,ZhouBatchelor1997}.
The displacement of one pattern is
\[
 d=a(h-v)+bh,
\]
and the cylinder identifies points differing by $md$. The marked edges
cross the cut at the bottom ports of the two marked sites in a fixed row.

\begin{proposition}[Physical meaning of the two-mark scalar]
\label{prop:two-mark-positive}
For this finite-width cylinder,
\begin{equation}
 \mathscr H=
 \sum_{\text{cylinder polygons using both marked edges}}
                              x_*^{|\gamma|}.                         \label{c4:Geom-2}
\end{equation}
The sum includes both contractible and winding polygons, with each
unoriented polygon counted once.
\end{proposition}

\begin{proof}
We first identify the two vacuum factors with physical half-cylinder
states. With twist $\eta=-i$, the no-defect matching transfer is
nonnegative and kills all closed loops. Its powers on the empty state
increase, as for the straight cut. They are bounded as well: each
remaining arc lifts to a planar arch below the zigzag cut and can be
closed by a fixed cap on the other side. Such a cap can be drawn in one
diagram row because bottom-to-bottom turns and horizontal passes have
positive weights for both column types. It joins distinct lifted stubs,
so closing gives a simple polygon through a fixed edge. The total
weight of these polygons is finite by \eqref{c4:Tail}. At fixed width
there are only finitely many endpoint choices for each matching; their
fixed cap lengths therefore bound its coefficient.

Every matching is reachable, by building its arcs from longer to shorter
intervals, using positive top-to-top turns and straight passes. The
rank-one convergence argument in Section~\ref{c4:cylinder} now applies.
Its normalized fixed vector is $\Psi(z)$: the polynomial fixed-vector
identity supplies it, and \eqref{c4:Pair-Pf} ensures regularity at these
sites. For inverse sites and auxiliary $Q^{-2}$ the two column types
are interchanged. The crossed transpose identity on charge zero thus
puts the same physical rows above the cut, with inverse twist and limit
$\Psi(z^{-1})^t$.

It remains to compute the phase of each finite-height diagram.
For an arc with endpoints $j<l$, in either half, its factor is
$q^{s_j}$ with $s_l=-s_j$. For a lower wrapping arc this follows from
\[
 q^{-s_j}\eta^{-s_j}=q^{s_j}.
\]
For upper arcs, use upward-positive spins and inverse twist: traversal
from the left endpoint at positive spin has turn $-2$ in square-diagram
quarter-turn units, with a rightward seam crossing for a wrap. This
gives the same endpoint rule. The pairings are noncrossing in cyclic,
hence chord, order. We may therefore redraw them as semicircles below
and above a horizontal line, with $P$ the leftmost slot. This redraw
computes phases separately for each original diagram. It preserves
connectivity; distinct cylinder diagrams retain their own physical
weights and are not identified, even if their redrawings coincide.

The two half-diagrams have matching spins except at $P,T$, where they
flip. A cycle cannot contain just one flip, so the two marks belong to
one alternating loop. Their prescribed signs orient its two paths
uniquely from $T$ to $P$, with opposite vertical arrival tangents at $P$.
Every other component is a freely oriented Jordan loop disjoint from
the marked component. It cannot separate $P$ from $T$, so its algebraic
intersection with the interval $P<j<T$ is zero. Thus it contributes zero
to $S_L$ in either orientation. Its two turn phases are $-i$ and $i$,
and cancel. Only the marked loop can survive.

For completeness, we compute its phase using the two oriented paths.
If $\gamma$ is either path, rounded to have a continuous unit tangent,
its signed turn satisfies
\[
 W(\gamma)=\Delta\arg(\gamma(t)-T)
                +\Delta\arg(P-\gamma(t)),
\]
with tangent limits at the endpoints. To obtain equality in lifted
arguments, consider the normalized secant
\[
 \frac{\gamma(t)-\gamma(s)}{|\gamma(t)-\gamma(s)|},
                    \qquad 0\le s<t\le1.
\]
Simplicity keeps the off-diagonal secants nonzero, and the regular
rounding gives a continuous extension to the closed ordered-time
triangle, with the unit tangent on the diagonal. A continuous argument
lift on that triangle compares the diagonal with its other two sides, giving
the displayed identity.

Both paths stay strictly to the right of $P$ before arrival. Their
$P-\gamma$ argument changes therefore sum to zero, since the arrival
tangents are opposite. For $\gamma-T$, cut the argument on the negative
horizontal axis. The two initial vertical arguments are opposite, as
are the limiting arguments at $P-T$ from the two half-planes. Every
intermediate cut crossing occurs at a slot $P<j<T$ and contributes
$-2\pi s_j$ to the lifted change. Consequently
\[
 W(\gamma_1)+W(\gamma_2)=-2\pi S_L.
\]
The arc phases have product $\exp(-i(W(\gamma_1)+W(\gamma_2))/4)
=\eta^{-S_L}$, canceling the insertion $\eta^{S_L}$. The marked loop
therefore has coefficient one. Taking the increasing finite-height
limits proves \eqref{c4:Geom-2} and its convergence.
\end{proof}

We now apply Theorem~\ref{c4:thm-two-mark} to these physical arrays.
Let $\iota=1$ if the first pattern site is unshifted and $\iota=0$ if it is
shifted. Omitting the marked sites gives the exact sector counts
\[
 \begin{aligned}
 (a_2,d_2)&=(a-\iota,b-1+\iota),\\
 (a_1,d_1)&=((m-1)a-\iota,(m-1)b-1+\iota).
 \end{aligned}
\]
Since $a\ge k/2$, we have $a_2\asymp k$, $a_1\asymp N$, and
$a_1\ge a_2$, with uniform comparison constants as $k\to\infty$.
Both ratios $d_i/a_i$ differ from $b/a\in[0,1]$ by $O(1/k)$,
uniformly including $b=1$. Hence, for fixed $C_0$ and $N\le k^{C_0}$,
\begin{equation}
                        |\mathscr H|\le k^{-4/3+o(1)}.                \label{c4:Geom-1}
\end{equation}
The exponent slack is uniform over these patterns: a failure for any
fixed positive slack at arbitrarily large $k$ would give a sequence
satisfying every hypothesis of Theorem~\ref{c4:thm-two-mark} and
contradict its conclusion.

\paragraph{Summing the marked displacements.}
Return to the plane translation classes used in \eqref{c4:Tail}.  Write
\(\ell(\gamma)=|\gamma|\) and \(w(\gamma)=x_*^{\ell(\gamma)}\).
\begin{proposition}[Polygon second length moment]\label{c4:prop-polygon-second-moment}
For unoriented simple polygons of the unit-edge honeycomb lattice,
counted modulo translations by \(\mathbb Z h+\mathbb Z v\),
\begin{equation}
                   \sum_{\operatorname{diam}\gamma\le r}|\gamma|^2 w(\gamma)
                             \le r^{2/3+o(1)}.                         \label{c4:Geom-3}
\end{equation}
as \(r\to\infty\), with Euclidean diameter and critical weight
\(w(\gamma)=x_*^{|\gamma|}\).
\end{proposition}

\begin{proof}
We first restrict to a diameter shell
$r/2\le\operatorname{diam}\gamma\le r$. On one fixed sublattice, a
polygon has $\ell(\gamma)/2$ vertices. Partition them by their missing
incident edge, which has three possible types, and let $m$ be the size
of a largest class. Then $m\ge\ell(\gamma)/6$. Let
$\mathcal P(\gamma)$ be the set of all ordered pairs on the chosen
sublattice with the same missing edge type, including diagonal pairs.
It contains the $m^2$ ordered pairs from the largest class, so
\[
             \ell(\gamma)^2\le36m^2\le36\lvert\mathcal P(\gamma)\rvert.
\]
We bound this last pair count by separating short, collinear, and
remaining displacements. Diagonal pairs belong to the short class.

For each first vertex, there are \(O(r)\) lattice vertices at distance
less than \(\sqrt r\).  There are also \(O(r)\) vertices within distance
\(r\) on the three triangular-grid lines through that vertex.  Thus the
pairs with short or collinear displacement contribute at most
\(C r\ell(\gamma)\) per polygon.  Their total shell contribution is
\[
 Cr\sum_{\operatorname{diam}\gamma\ge r/2}\ell(\gamma)w(\gamma)
       =r^{1/3+o(1)},
\]
by \eqref{c4:Tail}.  This is smaller than the required shell bound.

For every remaining pair, translate its first vertex to a fixed vertex
of the chosen sublattice and denote its displacement by \(d\).  A
honeycomb symmetry, followed if necessary by reflection, puts it in the
form
\[
 d=a(h-v)+bh,\qquad a,b\in\mathbb Z_{>0},\qquad b\le a.
\]
The strict positivity follows from having removed the three boundary
directions.  Set \(k=a+b\); then \(k\asymp|d|\), uniformly in the
direction.  The two vertices have the same missing edge type, so they
also share both occupied edge types.  Of these two types, at least one
is perpendicular to one of the two side steps \(h-v,h\).  Sum over the
bounded number of such choices, and start a pattern with that step.
Use \(a\) copies of \(h-v\) and \(b\) copies of \(h\) in the pattern.
Place its first port at the midpoint of the chosen occupied edge at the
first vertex.  Repetition puts the corresponding port at displacement
\(d\), in slot \(k+1\).  Both marked sites have the same type, as required
by the two-mark estimate.

Choose an even integer \(m\ge3\) so that \(N=mk\asymp r\) and
\(|md|>2r\).  Since \(k\gtrsim\sqrt r\), this choice has
\(3k\le N\le k^{C_0}\) for a fixed \(C_0\), after enlarging it to absorb
bounded constants.  Every counted polygon has diameter at most \(r\),
so its projection to this cylinder is injective.  Fixing the first marked
edge recovers its unique plane lift.  In particular, the anchored
polygons at displacement \(d\) count exactly the ordered vertex-pair
choices over translation classes; no factor proportional to the period
is introduced or lost.

Equations~\eqref{c4:Geom-1} and \eqref{c4:Geom-2} consequently bound the
mass for this displacement and edge choice by
\(Ck^{-4/3+o(1)}\), uniformly under the polynomial-width restriction.
For any fixed \(\varepsilon>0\), absorb the uniform exponent loss into
\(k^{\varepsilon}\).  There are \(O(t^2)\) lattice displacements in each
dyadic annulus \(t\le |d|<2t\), and hence
\[
 \sum_{\sqrt r\lesssim |d|\le r}|d|^{-4/3+\varepsilon}
   \le C_\varepsilon
      \sum_{\sqrt r/ C\le 2^j\le 2r}
            (2^j)^{2/3+\varepsilon}
   \le C_\varepsilon r^{2/3+\varepsilon}.
\]
Together with the exceptional-pair estimate this proves the same bound
for the shell.  Sum the shells below \(r\), including the finite
small-diameter contribution.  The resulting geometric series has the
same exponent.  Since \(\varepsilon>0\) was arbitrary, this proves
\eqref{c4:Geom-3}.
\end{proof}
